% TOP_PROBLEM: {"problemId": "problem.conformal-invariance-of-the-critical-random-cluster-model-with-q-4", "canonicalTitle": "Conformal Invariance of the Critical Random-Cluster Model with q <= 4", "releaseRank": 173, "primaryDomain": "probability_ergodic_dynamics", "primaryDomainLabel": "Probability, ergodic theory & dynamics", "displayStatus": "open_with_solved_subcases", "releaseStatus": "open_with_solved_subcases"}
% !TeX program = lualatex
% Definition and sources only; this file is not an LLM solution attempt.
\documentclass[11pt]{article}
\usepackage[margin=25mm]{geometry}
\usepackage{fontspec}
\setmainfont{Latin Modern Roman}
\usepackage{amsmath,amssymb,mathtools}
\usepackage{unicode-math}
\setmathfont{Latin Modern Math}
\usepackage{xurl,hyperref}
\hypersetup{colorlinks=true,urlcolor=blue}
\setcounter{secnumdepth}{0}
\setlength{\parindent}{0pt}
\setlength{\parskip}{5pt}
\setlength{\emergencystretch}{3em}
\begin{document}
\section*{\texorpdfstring{Conformal Invariance of the Critical Random-Cluster Model with q <= 4}{Conformal Invariance of the Critical Random-Cluster Model with q <= 4}}\label{173}
\subsection{Definitions and mathematical statement}
On a finite planar graph with boundary condition $\xi$, the FK measure gives a configuration $\omega\subseteq E$ weight proportional to
$p^{|\omega|}(1-p)^{|E|-|\omega|}q^{k_\xi(\omega)}$, where $k_\xi$ counts open clusters after boundary identifications. In a simply connected lattice approximation with marked boundary points, Dobrushin conditions produce an interface between wired primal and dual arcs. The intended statement is convergence of that interface at criticality to chordal $\mathrm{SLE}_{\kappa(q)}$ under mesh refinement. SLE is the Loewner chain driven by $\sqrt\kappa B_t$, and the customary FK relation is $\sqrt q=-2\cos(4\pi/\kappa)$, $4\le\kappa<8$, for $0<q\le4$; see \href{https://arxiv.org/html/1109.1549v4}{[E1]}. The catalogue omits the lower range of $q$, lattice, boundary setup, and curve topology. These must be taken from the cited formulation; the formula does not define a probability model for arbitrary negative $q$.
\par\smallskip\textit{Formulation and concept references:} \href{https://ar5iv.labs.arxiv.org/html/1707.00520}{[S1]}, \href{https://arxiv.org/html/1109.1549v4}{[E1]}.

\subsection{Short English statement}
Do critical planar random-cluster interfaces approach the conformally invariant SLE curve predicted by their cluster-weight parameter, with the model and boundary conventions properly fixed?
\subsection{Sources}
\begin{itemize}
\item[S1]H. Duminil-Copin, 'Lectures on the Ising and Potts models on the hypercubic lattice', arXiv:1707.00520 (2017), in Random Graphs, Phase Transitions and the Gaussian Free Field, Springer PIMS 123 (2020).
\url{https://ar5iv.labs.arxiv.org/html/1707.00520}.
\textit{Formulation source.}
\item[E1]Hugo Duminil-Copin and Stanislav Smirnov, Conformal invariance of lattice models, arXiv:1109.1549v4, FK interface and SLE discussion.
\url{https://arxiv.org/html/1109.1549v4}.
\textit{Added primary source: the random-cluster interface setup and parameter conventions. Consulted 24 September 2026.}
\end{itemize}
\end{document}
